\documentclass[11pt,a4paper]{article}
\usepackage{fontspec}
\usepackage{kotex}
\usepackage{geometry}
\usepackage{graphicx}
\usepackage{xcolor}
\usepackage{tcolorbox}
\usepackage{booktabs}
\usepackage{array}
\usepackage{longtable}
\usepackage{colortbl}
\usepackage{fancyhdr}
\usepackage{titlesec}
\usepackage{hyperref}
\usepackage{emoji}
\usepackage{amsmath}
\usepackage{amssymb}
\usepackage{enumitem}

\geometry{left=25mm, right=25mm, top=30mm, bottom=30mm}

\setmainfont{Noto Sans CJK KR}
\setsansfont{Noto Sans CJK KR}
\setmonofont{Noto Sans CJK KR}
\setemojifont{Noto Color Emoji}

\definecolor{primary}{HTML}{1976D2}
\definecolor{darkbrown}{HTML}{3E2723}
\definecolor{mediumbrown}{HTML}{5D4037}
\definecolor{lightbrown}{HTML}{8D6E63}
\definecolor{cream}{HTML}{FFF3E0}
\definecolor{lightgray}{HTML}{F5F5F5}
\definecolor{tableheader}{HTML}{E8E8E8}

\titleformat{\section}
  {\Large\bfseries\color{primary}}
  {\thesection.}{0.5em}{}
\titleformat{\subsection}
  {\large\bfseries\color{darkbrown}}
  {\thesubsection}{0.5em}{}
\titleformat{\subsubsection}
  {\normalsize\bfseries\color{mediumbrown}}
  {\thesubsubsection}{0.5em}{}

\renewcommand{\contentsname}{목차}
\renewcommand{\figurename}{그림}
\renewcommand{\tablename}{표}

\tcbset{
  tipbox/.style={colback=cream,colframe=primary,boxrule=1.5pt,arc=3mm,
    left=5mm,right=5mm,top=3mm,bottom=3mm,fonttitle=\bfseries\color{primary}},
  formulabox/.style={colback=lightgray,colframe=lightgray,boxrule=0pt,arc=2mm,
    left=5mm,right=5mm,top=3mm,bottom=3mm},
  summarybox/.style={colback=white,colframe=primary,boxrule=2pt,arc=4mm,
    left=5mm,right=5mm,top=5mm,bottom=5mm}
}

\hypersetup{colorlinks=true,linkcolor=primary,urlcolor=primary,bookmarks=true,unicode=true}


\pagestyle{fancy}
\fancyhf{}
\fancyhead[R]{\small\color{lightbrown}마이크로비트 Grove I2C 연결도 활용가이드 v1.0.0.0}
\fancyfoot[C]{\thepage}
\renewcommand{\headrulewidth}{0.4pt}

\begin{document}

\begin{titlepage}
\centering
\vspace*{3cm}
{\Huge\bfseries\color{primary}\emoji{electric-plug} 마이크로비트 Grove I2C 연결도}\\[0.5cm]
{\Huge\bfseries\color{primary} 활용가이드}
\\[1.5cm]
{\Large 시뮬레이터 버전: v1.0.3.0}
\\[2cm]
\begin{tcolorbox}[summarybox]
\centering
{\large 이 활용가이드는 마이크로비트 Grove I2C 연결도의 수업 적용 방법을 안내합니다.}
\end{tcolorbox}
\vfill
{\large 사물인터넷과 빅데이터}\\[0.3cm]
{\large 청주교육대학교 컴퓨터교육과}
\vspace{1cm}
\end{titlepage}

\section*{1. 수업 개요}

\begin{center}
\renewcommand{\arraystretch}{1.4}
\begin{tabular}{|p{2.8cm}|p{12.2cm}|}
\hline
\rowcolor{tableheader}
\textbf{항목} & \textbf{내용} \\
\hline
강좌 & 사물인터넷과 빅데이터 \\
\hline
적용 주차 & 2\textasciitilde{}4주차 \\
\hline
자료 역할 & 참고자료 \\
\hline
대상 & 4학년 (프로그래밍 비전공, 초등 예비교사) \\
\hline
수업 시간 & 45분 (1차시) \\
\hline
\end{tabular}
\end{center}

\subsection*{1.1 학습 목표}

$\bullet$ I2C 버스의 2선 구조(SDA, SCL)를 설명할 수 있다

$\bullet$ 센서별 I2C 주소의 역할을 이해할 수 있다

$\bullet$ Grove Shield의 물리적 연결 구조를 다이어그램으로 읽을 수 있다



\section*{2. 수업 시나리오 (45분)}

\begin{center}
\renewcommand{\arraystretch}{1.4}
\begin{tabular}{|p{1.5cm}|p{1.5cm}|p{9.1cm}|p{2.9cm}|}
\hline
\rowcolor{tableheader}
\textbf{단계} & \textbf{시간} & \textbf{활동} & \textbf{시뮬레이터 활용} \\
\hline
도입 & 5분 & "USB 케이블 1개 vs 센서마다 별도 케이블?" I2C 장점 소개 & — \\
\hline
관찰 & 10분 & 다이어그램에서 SDA/SCL 배선을 눈으로 추적 & 배선 구조 \\
\hline
학습 & 10분 & 각 센서 카드의 I2C 주소·데이터 항목 읽기 & 센서 카드 \\
\hline
활동 & 15분 & "주소가 같은 센서 2개를 연결하면?" 주소 충돌 토론 & — \\
\hline
정리 & 5분 & I2C 버스의 핵심 원리 요약, 실습 센서 배분 예고 & — \\
\hline
\end{tabular}
\end{center}


\section*{3. 학습 질문}

\subsection*{초급 (이해)}

$\bullet$ Q: I2C 버스에서 SDA와 SCL의 역할은 각각 무엇인가요?

$\bullet$ Q: I2C 주소는 왜 필요한가요?


\subsection*{중급 (적용)}

$\bullet$ Q: 6개 센서가 2개 선만으로 통신할 수 있는 이유를 설명해보세요.

$\bullet$ Q: 센서 카드에서 I2C 주소와 측정 데이터를 읽어보세요.


\subsection*{고급 (분석)}

$\bullet$ Q: 같은 I2C 주소를 가진 센서 2개를 동시에 연결하면 어떤 일이 벌어질까요?

$\bullet$ Q: I2C 방식과 개별 배선 방식의 장단점을 비교해보세요.


\subsection*{확장 (창의)}

$\bullet$ Q: I2C 이외에 센서를 연결하는 다른 통신 방식에는 무엇이 있을까요?

$\bullet$ Q: 버스 방식의 통신을 일상생활에서 찾아볼 수 있나요?



\section*{4. 학생 활동지}

\subsection*{활동: 마이크로비트와 Grove Shield의 I2C 센서 물리적 연결 구조 탐색하기}

\textbf{목표}: 마이크로비트와 Grove Shield의 I2C 센서 물리적 연결 구조를 보여주는 정적 다이어그램


\begin{center}
\renewcommand{\arraystretch}{1.4}
\begin{tabular}{|p{1.5cm}|p{5.6cm}|p{7.9cm}|}
\hline
\rowcolor{tableheader}
\textbf{단계} & \textbf{활동} & \textbf{기록} \\
\hline
1 & 화면 전체 구조 파악 & 주요 영역: \_\_\_\_\_\_\_\_\_\_\_\_\_\_\_\_\_ \\
\hline
2 & 핵심 기능 탐색 & 발견한 것: \_\_\_\_\_\_\_\_\_\_\_\_\_\_\_\_\_ \\
\hline
3 & 프로젝트와 연결 & 활용 방안: \_\_\_\_\_\_\_\_\_\_\_\_\_\_\_\_\_ \\
\hline
\end{tabular}
\end{center}


\section*{5. 평가 관찰 체크리스트}

\begin{center}
\renewcommand{\arraystretch}{1.4}
\begin{tabular}{|p{3.0cm}|p{6.3cm}|p{3.4cm}|p{2.3cm}|}
\hline
\rowcolor{tableheader}
\textbf{평가 항목} & \textbf{상} & \textbf{중} & \textbf{하} \\
\hline
핵심 개념 이해 & 관련 원리를 정확히 설명 & 대략적으로 이해 & 이해 부족 \\
\hline
실습 참여도 & 적극적으로 탐색·기록 & 지시에 따라 수행 & 소극적 참여 \\
\hline
프로젝트 연결 & 자기 프로젝트에 적용 방안 제시 & 일부 연결 시도 & 연결 미시도 \\
\hline
\end{tabular}
\end{center}


\section*{6. 수업 팁}

$\bullet$ \emoji{light-bulb} 실제 Grove 케이블·Shield가 있다면 다이어그램과 대조하며 보여주세요

$\bullet$ \emoji{light-bulb} SDA/SCL 배선을 손가락으로 따라가게 하면 이해가 빠릅니다

$\bullet$ \emoji{light-bulb} "왜 선이 2개면 충분한가?"라는 질문으로 주소 개념을 자연스럽게 도입하세요



\section*{7. 관련 자료}

\begin{center}
\renewcommand{\arraystretch}{1.4}
\begin{tabular}{|p{1.5cm}|p{7.3cm}|p{6.2cm}|}
\hline
\rowcolor{tableheader}
\textbf{\#} & \textbf{자료명} & \textbf{비고} \\
\hline
1 & 마이크로비트GroveI2C연결도\_퀵가이드 & 소프트웨어 사용 중 빠른 참조 \\
\hline
2 & 마이크로비트GroveI2C연결도\_사용설명서 & 전체 기능 상세 \\
\hline
3 & 마이크로비트GroveI2C연결도\_이론해설서 & 배경 이론 \\
\hline
4 & 마이크로비트GroveI2C연결도\_개발참조 & 기술 사양 \\
\hline
\end{tabular}
\end{center}


\vfill
\begin{center}
\small\color{lightbrown}
\textit{마이크로비트 Grove I2C 연결도 활용가이드}
\end{center}
\end{document}